% GENERATED FILE. Edit canonical lessons, then run npm run generate:books.
% canonical-source-hash: fnv1a64:48b23aed7b807e1b
% canonical-lessons: ZH-C363-zunzhong, ZH-C363-xiaosheng, ZH-C363-falu, ZH-C363-fanfa, ZH-C363-suzhi, ZH-R363-review-358-363

\chapter{Out in Public}
\label{ch:zh-out-in-public}
\hlchaptermodality{363}

\begin{chapteropening}
\textbf{By the end of this chapter:} \emph{I can talk about respect, speaking quietly, the law and how people behave.}

Complete the last lesson of chapter 363: I can talk about respect, speaking quietly, the law and how people behave.
\end{chapteropening}

\section[zūnzhòng]{zūnzhòng — to show respect for; to respect}

\label{lesson:ZH-C363-zunzhong}



\begin{quote}
Before the new one: say the Mandarin for impartial, then the Mandarin for a duty you owe.
\end{quote}

\subsection*{You'll want to know: zūnzhòng}
\emph{zūnzhòng} — \textquotedblleft{}to show respect for; to respect\textquotedblright{}.

\begin{quote}
\emph{wǒmen yīnggāi zūnzhòng biérén de xiǎngfǎ} — We should respect other people's ideas.
\end{quote}

Its characters come later; for now it is a word for the ear.

\begin{cousinweb}
\emph{zūnzhòng} is built from \emph{zūn} and \zh{重} 'heavy', both inherited from Middle Chinese, \emph{tswon} and \emph{drjowngX} (Guangyun, in Baxter's transcription), which glosses the first 'honoured, as against humble' and the second 'much' and 'careful'.
\end{cousinweb}

\begin{grammarlens}[title={What you've built}]
One more word for giving reasons.
\end{grammarlens}

\subsection*{Your turn}
\begin{itemize}
\raggedright
  \item \emph{Say it:} \emph{zūnzhòng}
  \item \emph{Say it:} \emph{zūnzhòng}, once more
  \item \emph{Say it:} \emph{wǒmen yīnggāi zūnzhòng biérén de xiǎngfǎ} — We should respect other people's ideas.
  \item \emph{Recall:} say the Mandarin for impartial, then the Mandarin for a duty you owe, then say \emph{zūnzhòng} again
\end{itemize}

\begin{tcolorbox}[breakable,colback=teal!4,colframe=teal!35!black,title={Before you move on}]
What does zūnzhòng mean? (\textquotedblleft{}To show respect for; to respect\textquotedblright{}.) Say it once more.
\end{tcolorbox}
\section[xiǎoshēng]{xiǎoshēng — in a low voice; quietly}

\label{lesson:ZH-C363-xiaosheng}



\begin{quote}
Before the new one: say the Mandarin for a duty you owe, then the Mandarin for to show respect for.
\end{quote}

\subsection*{You'll want to know: xiǎoshēng}
\emph{xiǎoshēng} — \textquotedblleft{}in a low voice; quietly\textquotedblright{}.

\begin{quote}
\emph{zài yīyuàn lǐ yīnggāi xiǎoshēng shuōhuà} — You should speak quietly in a hospital.
\end{quote}

Its characters come later; for now it is a word for the ear.

\begin{cousinweb}
\emph{xiǎoshēng} is built from \zh{小} 'small' and \emph{shēng}, both inherited from Middle Chinese, \emph{sjewX} and \emph{syeng} (Guangyun, in Baxter's transcription), which glosses the first 'tiny' and the second 'a sound'.
\end{cousinweb}

\begin{grammarlens}[title={What you've built}]
One more word for giving reasons.
\end{grammarlens}

\subsection*{Your turn}
\begin{itemize}
\raggedright
  \item \emph{Say it:} \emph{xiǎoshēng}
  \item \emph{Say it:} \emph{xiǎoshēng}, once more
  \item \emph{Say it:} \emph{zài yīyuàn lǐ yīnggāi xiǎoshēng shuōhuà} — You should speak quietly in a hospital.
  \item \emph{Recall:} say the Mandarin for a duty you owe, then the Mandarin for to show respect for, then say \emph{xiǎoshēng} again
\end{itemize}

\begin{tcolorbox}[breakable,colback=teal!4,colframe=teal!35!black,title={Before you move on}]
What does xiǎoshēng mean? (\textquotedblleft{}In a low voice; quietly\textquotedblright{}.) Say it once more.
\end{tcolorbox}
\section[fǎlǜ]{fǎlǜ — law; the law}

\label{lesson:ZH-C363-falu}



\begin{quote}
Before the new one: say the Mandarin for to show respect for, then the Mandarin for in a low voice.
\end{quote}

\subsection*{You'll want to know: fǎlǜ}
\emph{fǎlǜ} — \textquotedblleft{}law; the law\textquotedblright{}.

\begin{quote}
\emph{wǒmen bìxū zūnshǒu fǎlǜ} — We must obey the law.
\end{quote}

Its characters come later; for now it is a word for the ear.

\begin{cousinweb}
\emph{fǎlǜ} is built from \zh{法} 'a rule' and \emph{lǜ}, both inherited from Middle Chinese, \emph{pjop} and \emph{lwit} (Guangyun, in Baxter's transcription), which glosses the first 'a rule' and the second 'the pitch pipes' and 'a law'.
\end{cousinweb}

\begin{grammarlens}[title={What you've built}]
One more word for giving reasons.
\end{grammarlens}

\subsection*{Your turn}
\begin{itemize}
\raggedright
  \item \emph{Say it:} \emph{fǎlǜ}
  \item \emph{Say it:} \emph{fǎlǜ}, once more
  \item \emph{Say it:} \emph{wǒmen bìxū zūnshǒu fǎlǜ} — We must obey the law.
  \item \emph{Recall:} say the Mandarin for to show respect for, then the Mandarin for in a low voice, then say \emph{fǎlǜ} again
\end{itemize}

\begin{tcolorbox}[breakable,colback=teal!4,colframe=teal!35!black,title={Before you move on}]
What does fǎlǜ mean? (\textquotedblleft{}Law; the law\textquotedblright{}.) Say it once more.
\end{tcolorbox}
\section[fànfǎ]{fànfǎ — to break the law}

\label{lesson:ZH-C363-fanfa}



\begin{quote}
Before the new one: say the Mandarin for in a low voice, then the Mandarin for the law.
\end{quote}

\subsection*{You'll want to know: fànfǎ}
\emph{fànfǎ} — \textquotedblleft{}to break the law\textquotedblright{}.

\begin{quote}
\emph{hējiǔ yǐhòu kāichē fànfǎ} — Driving after drinking is against the law.
\end{quote}

Its characters come later; for now it is a word for the ear.

\begin{cousinweb}
\emph{fànfǎ} is built from \emph{fàn} and \zh{法} 'a rule', both inherited from Middle Chinese, \emph{bjomX} and \emph{pjop} (Guangyun, in Baxter's transcription), which glosses the first 'to offend against' and 'to intrude', and the second 'a rule'.
\end{cousinweb}

\begin{grammarlens}[title={What you've built}]
One more word for giving reasons.
\end{grammarlens}

\subsection*{Your turn}
\begin{itemize}
\raggedright
  \item \emph{Say it:} \emph{fànfǎ}
  \item \emph{Say it:} \emph{fànfǎ}, once more
  \item \emph{Say it:} \emph{hējiǔ yǐhòu kāichē fànfǎ} — Driving after drinking is against the law.
  \item \emph{Recall:} say the Mandarin for in a low voice, then the Mandarin for the law, then say \emph{fànfǎ} again
\end{itemize}

\begin{tcolorbox}[breakable,colback=teal!4,colframe=teal!35!black,title={Before you move on}]
What does fànfǎ mean? (\textquotedblleft{}To break the law\textquotedblright{}.) Say it once more.
\end{tcolorbox}
\section[sùzhì]{sùzhì — a person's quality; upbringing}

\label{lesson:ZH-C363-suzhi}



\begin{quote}
Before the new one: say the Mandarin for the law, then the Mandarin for to break the law.
\end{quote}

\subsection*{You'll want to know: sùzhì}
\emph{sùzhì} — \textquotedblleft{}a person's quality; upbringing\textquotedblright{}.

\begin{quote}
\emph{wǒ juéde zhège rén sùzhì hěn gāo} — I think he's a very well brought-up person.
\end{quote}

Its characters come later; for now it is a word for the ear.

\begin{cousinweb}
\emph{sùzhì} is built from \emph{sù} and \emph{zhì}, both inherited from Middle Chinese, \emph{suH} and \emph{tsyit} (Guangyun, in Baxter's transcription), which glosses the first 'empty' and 'silk', and the second 'plain' and 'sincere'.
\end{cousinweb}

\begin{grammarlens}[title={What you've built}]
One more word for giving reasons.
\end{grammarlens}

\subsection*{Your turn}
\begin{itemize}
\raggedright
  \item \emph{Say it:} \emph{sùzhì}
  \item \emph{Say it:} \emph{sùzhì}, once more
  \item \emph{Say it:} \emph{wǒ juéde zhège rén sùzhì hěn gāo} — I think he's a very well brought-up person.
  \item \emph{Recall:} say the Mandarin for the law, then the Mandarin for to break the law, then say \emph{sùzhì} again
\end{itemize}

\begin{tcolorbox}[breakable,colback=teal!4,colframe=teal!35!black,title={Before you move on}]
What does sùzhì mean? (\textquotedblleft{}A person's quality; upbringing\textquotedblright{}.) Say it once more.
\end{tcolorbox}
\section[(dialogue)]{Looking back: from Real Manners to Out in Public}

\label{lesson:ZH-R363-review-358-363}



\begin{quote}
Say the Mandarin for good manners and genuine.
\end{quote}

\subsection*{Your turn}
\begin{itemize}
\raggedright
  \item \emph{Say it:} good manners, genuine, from the bottom of your heart, modest about yourself, face (your standing)
  \item \emph{Say it:} the rules of a game, to abide by, to break a rule, to put a limit on, morality
  \item \emph{Say it:} shared by everyone, necessary, there's no need to, the occasion you are in, to jump the queue
  \item \emph{Say it:} to say hello to someone, to look after your guests, to bow, to propose a toast, to look down on someone
  \item \emph{Say it:} in a loud voice, must (a firm requirement), in a mess, impartial, a duty you owe
  \item \emph{Say it:} to show respect for, in a low voice, the law, to break the law, a person's character and upbringing
  \item \emph{Read it:} \zh{真} \zh{必} \zh{须} — each one aloud, then a word that holds it
\end{itemize}

\begin{tcolorbox}[breakable,colback=teal!4,colframe=teal!35!black,title={Before you move on}]
Good manners? (\emph{lǐmào}.) Genuine? (\textbf{\zh{真正}}, \emph{zhēnzhèng}.) From the bottom of your heart? (\textbf{\zh{真心}}, \emph{zhēnxīn}.) Modest about yourself? (\emph{qiānxū}.) Face (your standing)? (\textbf{\zh{面子}}, \emph{miànzi}.) The rules of a game? (\emph{guīzé}.) To abide by? (\emph{zūnshǒu}.) To break a rule? (\emph{wéifǎn}.) To put a limit on? (\emph{xiànzhì}.) Morality? (\emph{dàodé}.) Shared by everyone? (\emph{gōnggòng}.) Necessary? (\textbf{\zh{必要}}, \emph{bìyào}.) There's no need to? (\textbf{\zh{不必}}, \emph{búbì}.) The occasion you are in? (\emph{chǎnghé}.) To jump the queue? (\emph{chāduì}.) To say hello to someone? (\emph{dǎ zhāohu}.) To look after your guests? (\emph{zhāodài}.) To bow? (\emph{jūgōng}.) To propose a toast? (\emph{jìngjiǔ}.) To look down on someone? (\textbf{\zh{看不起}}, \emph{kànbuqǐ}.) In a loud voice? (\emph{dàshēng}.) Must (a firm requirement)? (\textbf{\zh{必须}}, \emph{bìxū}.) In a mess? (\emph{luàn}.) Impartial? (\textbf{\zh{公正}}, \emph{gōngzhèng}.) A duty you owe? (\emph{yìwù}.) To show respect for? (\emph{zūnzhòng}.) In a low voice? (\emph{xiǎoshēng}.) The law? (\emph{fǎlǜ}.) To break the law? (\emph{fànfǎ}.) A person's character and upbringing? (\emph{sùzhì}.)
\end{tcolorbox}
